\section{Proofs for lossless bidirectional charging}\label{app:v2g}
The assumptions and result statements remain in Section~\ref{sec:v2g}.

\subsection{Why the projection is exact}
\begin{proof}[Proof of Lemma~\ref{lem:v2ggross}]
A paired bus has signed charge in $[a_i+b_j-B,a_i]$: early headroom
gives the upper bound, and enough inventory for B gives the lower bound.
An A-only bus has $q\in[a_i-B,a_i]$; a B-only bus has
$q\in[b_j-B,0]$. These endpoint conditions keep inventory in $[0,B]$
before and after the remaining service. A monotone action between those
inventories is also feasible throughout the early window. Terminal demand
is that bus's service consumption minus $q$, between zero and $B$.

All upper endpoints are nonnegative and sum to $P$. The sum of positive
lower endpoints is $R$, while the sum of all lower endpoints is
$T=S-N_mB$. Thus total possible negative excursion is $R-T$. For
$E\ge R$, set negative coordinates to zero and distribute purchases
among the independent nonnegative intervals, whose sum spans $[R,P]$.
Gross energy is then $E$. For $T\le E<R$, keep mandatory positive
coordinates at their lower endpoints and distribute discharge of magnitude
$R-E$ among the remaining negative intervals. This is possible exactly
because $R-E\le R-T$ and has gross energy $2R-E$.

Conversely, total positive purchases are at least $R$, so any allocation
with net $E$ has gross energy at least $2R-E$ and at least $E$.
The constructions attain both regimes in \eqref{eq:v2ggross}. Additional
charge/discharge cycles cannot reduce gross energy and are unnecessary.
\end{proof}

\begin{proof}[Proof of Theorem~\ref{thm:v2g}]
The lemma and early gross budget give $T\le E\le P$, $E\le K$
and $E\ge2R-K$. All buses are present for the full early window;
their monotone actions serialize in at most one hour when $K>0$.
For $K=0$, the inequalities force $R=0$, $E=0$ and zero gross energy,
without division. Every bus's terminal refill is nonnegative and at most
$B$. They serialize at $H$ exactly when $S-E\le H$; at $H=0$ no
positive terminal demand can be met. With a nonempty positive B batch,
$E\le P<S$ makes that zero-terminal-power case infeasible.

For a fixed matching, the exact interval is consequently
\[
 [\max\{S-H,S-N_mB,2R(M)-K\},U].
\]
Only $R(M)$ varies at fixed cardinality. Its attained minimum $r_m$
gives the widest interval; all such intervals have common upper endpoint
$U$. Their union is \eqref{eq:v2gprojection}. Since $r_m$ is
nondecreasing, its threshold minimum is $t_0$; feasibility is precisely
$t_0\le U$, attained by a standalone fleet when needed.

Standard allowed-edge bipartite matching convexity gives discretely convex
$r_m$. Interpolate it linearly on $[0,\nu]$. Taking the maximum of
the resulting convex function $2r-K$, the affine function $S-(2n-m)B$
and the constant $S-H$ gives a convex function whose integer values are
$t_m$. Each component is nondecreasing, so the thresholds are also
nondecreasing. Zero increments can occur only in an initial block.
After retaining its cheapest cardinality, the positive threshold increments
are nondecreasing. Frontier slopes $-f/(t_m-t_{m-1})$ are therefore
nondecreasing, followed by zero. Each endpoint is physical and each edge
is its whole-plan mixture; every physical branch lies above this convex
graph. This proves the lower-boundary assertion without introducing
nonexistent or resource-infeasible cardinalities.
\end{proof}

\begin{proof}[Proof of Corollary~\ref{cor:v2gsupport}]
Strict convexity of $G$ gives one minimizing hull load. At a retained
threshold or on the final horizontal segment, a physical plan attains the
frontier. In an adjacent open threshold interval, $f>0$ makes its
interpolated cost strictly below every physical branch at that load.
The compact finite union of physical intervals then has a strictly larger
minimum. The classical support criterion applies only with the stated
signed-domain and price qualifications. At fee zero the empty matching
physically attains the entire interval at zero operating cost.
\end{proof}

\subsection{Identical paid idle buses}
\begin{proof}[Proof of Proposition~\ref{prop:v2gidle}]
Consider a matching of size $m$ with $k$ idle buses. Its total fleet is
$N=2n-m+k$, so the cap implies $k\le m$. The box argument of
Lemma~\ref{lem:v2ggross} still applies: idle buses add only negative
headroom. The matching's interval is therefore
\[
 [\max\{S-H,S-NB,2R-K\},U],
\]
at operating cost $fN$. Delete any $k$ matching edges and use the resulting
$m-k$ matching without idle buses. This fleet has the same size $N$ and
fee. Deleting edges cannot increase $R$, while the inventory term $S-NB$
and upper endpoint $U$ stay fixed. Its feasible interval contains the
original interval. Thus every projected point with idle buses has an
identical load and cost without them. The reverse inclusion is immediate.
\end{proof}
